%%%%%%%%%%%%%%%%%%%%%%%%%%%%%%%%%%%%%%%%%%%%%%%%%%%%%
%			HLAVIČKA								%
%%%%%%%%%%%%%%%%%%%%%%%%%%%%%%%%%%%%%%%%%%%%%%%%%%%%%
\documentclass[openany,12pt]{memoir}

\usepackage[utf8]{inputenc} 
\usepackage[czech]{babel}
\usepackage[T1]{fontenc}
\usepackage[top=1.5cm, bottom=2cm, left=2cm, right=2cm]{geometry}  % --> NASTAVENÍ OKRAJŮ
\usepackage{fancyhdr}
\usepackage{graphicx}
\usepackage{lmodern}
\usepackage{xwatermark}
\usepackage{xcolor}
\usepackage{changepage}
\usepackage{newunicodechar}



%%%%%% Package na zpěvník
\usepackage[full]{leadsheets}%http://mirrors.nic.cz/tex-archive/macros/latex/contrib/leadsheets/leadsheets_en.pdf   --> dokumentace	
\definesongtitletemplate{empty}{} 
\setchords{
format = \bfseries,   %tučné akordy
minor = {mi},% 
input-notation = {german},%
output-notation = {german}%
}
\definesongtitletemplate{empty}{} 

\newlength{\drop}
\newwatermark[allpages,color=red!50,angle=0,scale=2, xpos=0,ypos=0]{\includegraphics[width=5cm]{obr/pozadi.jpg}} %--> dvojka na pozadí


%%%% Vlastní příkazy
\newcounter{Slokočet}   %Automatické číslování slok
\newcommand{\mezera}{\vspace*{0.5cm}}   %Horizontální odsazení slok
\newcommand{\stred}{5.2cm}   %%% Na zarovnání slok doprostřed, pozn. automatičtější zarovnávání na střed nejde
\newcommand{\refren}{\mezera \noindent \textbf{R:} } %refrén
\newcommand{\sloka}{\mezera \noindent \addtocounter{Slokočet}{1} \arabic{Slokočet}. } 	%sloka, která se automaticky čísluje
\newcommand{\ssloka}{\mezera \noindent}  % vlastní číslo sloky
\newcommand{\carka}{,\:}
\newcommand{\m}[1]{\color{white}{#1}}  %Pro akordy

%%%%%%%%%%%%%%%%%%%%%%%%%%%%%%%%%%%%%%%%%%%%%%%%%%%%%%
%			NÁVOD									 %
%%%%%%%%%%%%%%%%%%%%%%%%%%%%%%%%%%%%%%%%%%%%%%%%%%%%%%
%1. Věci v hlavičce IGNOROVAT
%2. Píseň psát do prostoru mezi \begin{song} a \end{song}
%3. další řádek se značí dvěma odsazeníma (= dvakrát stisknout enter)
%4. \refren vždy na začátku refrenu a \sloka na začátek sloky (automaticky se čísluje)
%  \ = alt gr + q ; [/] = alt gr f/g ; {/} = alt gr + b/n; ^ = alt gr + 3 + mezera
%Cokoliv napíšete do ^{  } se bude brát jako akord
%když se toto bude dotýkat nějakého slova (nebude mezi tím a slovem mezera)
%tak se akordy zjeví nad slovem, ale pište to před slovo
%Když se to nedotýká slova, tak akord lítá ve vzduchu a vytiskne se větší mezera
%První možnost je asi preferovanější
%5. Akordy stačí psát jen do první sloky, když se nezmění -- kytaristi to zvládnou



%<++>
%\usepackage{subfiles}
%</++>

\begin{document}

\pagestyle{empty}
%\addtocontents{toc}{\protect\thispagestyle{empty}}
%\tableofcontents \thispagestyle{empty}\newpage
\newgeometry{top=1.5cm, bottom = 0cm, left = 2cm, right = 2cm}




%Pro kompilaci je důležitý, aby ten input a \end{document} zůstali na posledních
%dvou řádcích!!!!
\begin{song}{title=\centering Let It Be \\\normalsize The Beatles  \vspace*{-0.3cm}}  %% sem se napíše jméno songu a autor
\moveright 4.8cm \vbox{      %Varianta č. 1  ---> Jeden sloupec zarovnaný na střed	

\sloka 
	When I ^{C\,\,}find myself in ^{G\,\,\,\,\,\,}times of trouble

	^{Ami\,\,\,\,}Mother Mary ^{Fmaj7\,\,}Comes to me
	
	^{C\,\,\,\,\,\,\,\,\,\,\,}Speaking words of ^{G\,\,\,\,\,\,\,\,\,\,}wisdom let it ^{F}be ^{C\,Dmi7\,C}
	
	^{C}And in my hour of ^{G\,\,\,\,\,\,\,\,\,\,}darkness

	She is ^{Ami\,\,\,\,}standing right in ^{Fmaj7}front of me

	^{C\,\,\,\,\,\,\,\,}Speaking words of ^{G\,\,\,\,\,\,\,\,\,\,}wisdom let it ^{F}be. ^{C\,Dmi7\,C}

\refren
	Let it ^{Ami}be let it ^{G}be let it ^{F}be let it ^{C}be

	^{C\,\,\,\,}Whisper words of ^{G\,\,\,\,\,\,\,\,\,\,}wisdom let it ^{F}be. ^{C\,Dmi7\,C}

\sloka
	And when the broken hearted people
   	
   	Living in the world agree
	
	There will be an answer let it be
   	
   	For though they may be parted
   	
   	There is still a chance that they will see
  	
  	There will be an answer let it be.

\moveright -0.2cm \vbox{\ssloka \textbf{R\textsubscript{2}:}
	Let it be let it be let it be let it be}
    
    There will be an answer let it be.

\refren


\refren

\sloka
	And when the night is cloudy
   	
   	There is still a light that shines on me
	
	Shine until the morrow let it be
   	
   	I wake up to the sound of music
   	
   	Mother Mary comes to me
   	
   	Speaking words of wisdom let it be.

\moveright -0.2cm \vbox{\ssloka \textbf{R\textsubscript{2}:}}


\moveright -0.2cm \vbox{\ssloka \textbf{R\textsubscript{2}:}}


}
\setcounter{Slokočet}{0}
\end{song}



\end{document}
